\section{Discussão e Análise de Resultados}
\label{cap:discussao}

...
\section{Discussão e Análise de Resultados}

A análise do cenário inicial, guiada estritamente pela Teoria da Probabilidade clássica, permitiu mapear as chances estáticas dos eventos antes de qualquer imprevisto. No experimento do lançamento de moedas, a probabilidade teórica projetou uma chance clara de 12,5\% para a ocorrência tripla, servindo como uma linha de base sólida para os cálculos. Da mesma forma, no cenário analítico de diagnósticos, a probabilidade inicial indicava que apenas 1\% da população possuía determinada característica, estabelecendo uma fundação matemática estática antes da inserção de novas variáveis empíricas.

Contudo, a introdução de novas evidências exigiu a aplicação do Teorema de Bayes, alterando drasticamente a percepção de risco mapeada inicialmente. Conforme ilustrado na Tabela \ref{tab:comparativo_bayes}, a probabilidade prévia de um indivíduo pertencer ao grupo de risco (1\%) sofreu um choque de realidade ao ser confrontada com um teste positivo, saltando para 16,10\%. Essa tabela comprova como a nova evidência calibra o sistema, demonstrando que confiar apenas na probabilidade estática inicial ou na alta precisão isolada de um teste (95\%) levaria a uma interpretação totalmente precipitada e arriscada.

\begin{table}[ht]
    \centering
    \caption{Comparação de Probabilidades: Cenário Estático vs. Atualização Bayesiana}
    \label{tab:comparativo_bayes}
    \begin{tabular}{|l|c|c|}
        \hline
        \textbf{Métrica Analisada} & \textbf{Probabilidade Antes} & \textbf{Probabilidade Depois (Bayes)} \\ \hline
        Prevalência Base na População & 1,00\% & - \\ \hline
        Positivo Verdadeiro (Sensibilidade) & - & 95,00\% \\ \hline
        Falso Positivo & - & 5,00\% \\ \hline
        \textbf{Chance Real Atualizada} & \textbf{1,00\%} & \textbf{16,10\%} \\ \hline
    \end{tabular}
\end{table}

Com as probabilidades atualizadas em mãos, os dados foram submetidos à análise de Valor Esperado para projetar tendências a longo prazo sob a ótica da Lei dos Grandes Números. Para evidenciar esse comportamento, elaborou-se um algoritmo para análise de dados realizando um milhão de simulações com pesos e valores variados, estruturado da seguinte forma:

[placeholder para o bloco de código Python do Valor Esperado]

Observando o gráfico gerado por essa simulação (Figura \ref{fig:valor_esperado}), conclui-se que a média dos elementos aleatórios e sua comparação com o valor esperado não são inicialmente determinísticos. Com poucos dados gerados, cada sorteio possui um impacto extremo, resultando em violentas oscilações. Porém, quanto mais iterações são realizadas (principalmente ultrapassando a abscissa $10^2$ e ordenada $10^4$), mais as curvas tendem a se estabilizar e se alinhar uma sobre a outra, possuindo uma diferença insignificante e provando que a expectativa de longo prazo respeita a matemática calculada.

\begin{figure}[ht]
    \centering
    [placeholder para a imagem do gráfico comparativo de Valor Esperado]
    \caption{Convergência do Valor Esperado ao longo das iterações logarítmicas.}
    \label{fig:valor_esperado}
\end{figure}

Por fim, a Variância atuou como o termômetro fundamental para medir a instabilidade da operação e atestar o nível de risco dos valores gerados. Para demonstrar a variabilidade de um grande volume de dados estocásticos, a variância empírica e a teórica foram calculadas cumulativamente a cada iteração, por meio do seguinte script:

[placeholder para o bloco de código Python da Variância]

Conforme atestado pelo decaimento exponencial na Figura \ref{fig:variancia}, ocorre uma primeira oscilação na variância que reflete a mesma instabilidade brutal visualizada inicialmente no gráfico de Valor Esperado. Contudo, com o passar da série de simulações, o alcance de variabilidade diminui e o gráfico segue caindo cada vez mais próximo de zero, sem nunca tocá-lo. Isso comprova de forma incontestável a nossa análise de risco: decisões estruturais baseadas em amostragens curtas sofrem de alta variância (alto risco), tornando-se seguras apenas quando o volume massivo de testes força a redução estrita do caos rumo à constância da média.

\begin{figure}[ht]
    \centering
    [placeholder para a imagem do gráfico de Variância da Média Amostral]
    \caption{Decaimento da Variância da Média em escala logarítmica.}
    \label{fig:variancia}
\end{figure}
...
